% =============================================================================
\section{Vision: A Research Roadmap}
\label{sec:vision}
% =============================================================================

The WRP matrix in \Cref{tab:interactions} reveals two structurally
sparse cells---Router $\times$ Pool and Pool $\times$ Router---and
several interactions in the dense cells that our publications have
identified but not yet exploited.  Below we propose twenty-one research
directions that follow directly from these gaps, each framed around
what the semantic router's signal architecture and the WRP
cross-dimensional view uniquely enable.

\Cref{tab:maturity} assigns each opportunity a maturity tier.
Throughout this section, projected benefits are estimates grounded
in our prior measurements unless a specific external citation is
given.

\begin{table}[t]
\centering
\caption{Maturity and feasibility of the twenty-one proposed
opportunities.  \textbf{Tier} indicates the primary barrier:
\emph{engineering} (building blocks exist, integration needed),
\emph{research} (open technical questions).}
\label{tab:maturity}
\small
\renewcommand{\arraystretch}{1.15}
\begin{tabularx}{\textwidth}{c l c X}
\toprule
\textbf{\#} & \textbf{Opportunity} & \textbf{Tier} & \textbf{Primary barrier} \\
\midrule
\rowcolor{routergreen!8}
2  & HaluGate model reputation routing    & Engineering & Sample size per (model, domain) cell \\
5  & Token-budget enforcement for agents   & Engineering & Domain-specific budget model calibration \\
\rowcolor{routergreen!8}
8  & Follow-the-sun pool rebalancing       & Engineering & Sliding-window CDF estimation + FleetOpt \\
12 & KV-cache retention directives         & Engineering & Retention-directive API integration \\
\rowcolor{routergreen!8}
15 & Governance-as-code                    & Engineering & Extend DSL with WRP constraints \\
19 & Request-level RBAC enforcement         & Engineering & Gateway hook + JWT integration \\
\midrule
\rowcolor{poolorange!8}
1  & Offline RL for session-length routing  & Research & Credit assignment + off-policy evaluation \\
3  & Cumulative multi-turn risk scoring     & Research & Threshold calibration on adversarial benchmarks \\
\rowcolor{poolorange!8}
4  & Failure-class-aware routing            & Research & Multi-dimensional failure weighting \\
\textbf{17} & \textbf{Fleet-learned joint tool--model optimization} & \textbf{Research} & \textbf{Joint (tool, model) outcome table sparsity} \\
\rowcolor{poolorange!8}
18 & Gateway-coordinated agent loops       & Research & End-to-end fleet integration unevaluated; subsystems validated separately \\
\rowcolor{poolorange!8}
20 & Fleet-wide caller reputation           & Research & Heterogeneous security-event scoring + decay calibration \\
21 & Router-enforced behavioral commitments & Research & Boundary-declaration NLU classifier + false-positive risk \\
\rowcolor{poolorange!8}
6  & Pool-state inference from observables  & Research & Prefill-rate regression + cold start \\
7  & Output-length-aware pool routing       & Research & Output-length prediction variance \\
\rowcolor{poolorange!8}
9  & Pool-aware cascading                   & Research & Cascade-abort threshold + grace period \\
10 & Online $(\gamma, \boundary^*)$ co-adaptation & Research & Oscillation control + validation gap \\
\rowcolor{poolorange!8}
11 & Mixed-archetype fleet provisioning     & Research & Multi-tenant FleetOpt extension \\
13 & Archetype-driven proactive scaling     & Research & Archetype-aware forecast vs.\ aggregate baseline \\
\rowcolor{poolorange!8}
14 & Energy as a routing objective          & Research & tok/W signal calibration per pool config \\
16 & Closed-loop self-adaptation            & Research & Knob interaction strength + composite reward \\
\bottomrule
\end{tabularx}
\end{table}


% -----------------------------------------------------------------------------
\subsection{Workload $\times$ Router}
% -----------------------------------------------------------------------------

\begin{opportunity}[Offline RL for session-length-minimizing routing]
\label{opp:session-cascade}
In agent workloads, a long session is often the consequence of poor
routing decisions at early turns.  If the router selects a weak
model at turn~2 and the model fails a tool call, the agent retries
at turn~3 with a longer prompt, which fails again, and so on.  A
session that could have completed in 3~turns with the right model
and tool selection at the outset may instead take 8~turns---each
additional turn growing the prompt
(ServeGen~\cite{servegen2025} reports 3--10$\times$ more LLM calls
per user request for agents than chatbots) and consuming
proportionally more compute.

The router observes the full trajectory of every completed session:
which model served each turn, which tools were invoked
(OATS~\cite{oats2026}), whether each tool call succeeded, and how
many turns the session took.  This is a natural offline RL dataset.
The state is (turn number, tools tried so far, cumulative context
length, domain); the action is (model, tool) selection; the reward
is negative session length (or task-completion / session-length for
quality-adjusted optimization).  The workload identity signals from
\Cref{sec:workload:identity} enrich this state space at zero
inference cost: the system-prompt fingerprint provides an instant
domain prior, the tool-definition catalog identifies the agent's
operational mode, and the conversation depth from the messages
array gives exact session maturity---all available before the
router inspects content.  The router's fleet-wide visibility
is the structural advantage over agent-side SDKs: it collects
trajectories from thousands of sessions across all tenants, giving
it the training volume that no single agent instance can match.

An offline RL policy (e.g., conservative Q-learning or decision
transformer) trained on these trajectories learns which early-turn
decisions lead to shorter sessions: ``for coding tasks, routing
turn~1 to the frontier model and selecting \texttt{semantic\_search}
over \texttt{grep} reduces median session length from 7 to
4~turns.''  The learned policy replaces OATS's per-turn greedy
selection with a \emph{per-stage} strategy that adapts model and
tool selection as the session evolves---the same progression
Aragog~\cite{aragog2025} demonstrates for throughput (50--217\%
improvement by decoupling one-time configuration from per-stage
scheduling), applied here to session-length minimization.

\begin{frisk}
The router already logs per-turn model and tool-call outcomes for
OATS and the Feedback Detector; assembling these into session
trajectories is bookkeeping.  Offline RL is well suited to this
setting: recent work shows that conservative Q-learning variants
learn effective scheduling policies from logged data
alone~\cite{offlinerl2025scheduling}, generalizing even from
suboptimal (random-heuristic) training trajectories---closely
matching the quality of router logs collected under a non-optimal
greedy policy.  The research questions are:
(a)~credit assignment---which of the 8~turns caused the extra
5~turns?  Offline RL handles this via temporal-difference learning,
but noisy agent environments may require variance reduction.
(b)~Off-policy evaluation---the learned policy must be validated
before deployment.  Importance-weighted estimators on held-out
sessions provide offline evaluation without live traffic risk.
The comparison baseline is OATS with single-turn greedy selection;
the metric is median session length (turns and tokens) for
completed tasks on SWE-bench and BFCL.
\end{frisk}
\end{opportunity}

\begin{opportunity}[HaluGate-driven model reputation routing]
\label{opp:halugate-reputation}
The Factcheck Classifier~\cite{factcheck2026} identifies
fact-seeking prompts at request time ($\sim$12\,ms).
HaluGate~\cite{halugate2025} validates responses at token level
(76--162\,ms overhead) and produces a per-response hallucination
verdict.  Today, these verdicts are logged but not fed back into
routing decisions.

The router processes every response through HaluGate for
fact-seeking traffic.  Over time, it accumulates (model, domain,
hallucinated) triples---enough to maintain a per-model,
per-domain accuracy estimate.  For example: ``Model~A hallucinates
on medical queries 12\% of the time but only 2\% on coding;
Model~B is the reverse.''  The opportunity is to use this estimate
as a live routing signal: for each fact-seeking request, the
router queries the per-model accuracy table for the request's
domain and routes to the model with the lowest hallucination rate,
subject to cost and latency constraints.

\begin{frisk}
The data collection is zero-overhead (HaluGate already runs).
The per-model accuracy table is a counter per (model, domain) pair,
updated on each verdict.  The engineering question is sample
size: how many verdicts per (model, domain) cell are needed before
the estimate is reliable?  A Bayesian approach (Beta prior,
Bernoulli likelihood) gives calibrated confidence intervals from
small samples.  The risk is distributional shift: model updates
invalidate historical accuracy estimates.  A sliding window
(e.g., last 1,000 verdicts per cell) handles this.
\end{frisk}
\end{opportunity}

\begin{opportunity}[Cumulative multi-turn risk scoring]
\label{opp:cumulative-risk}
SafetyL1~\cite{safetyl1_2026} and SafetyL2~\cite{safetyl2_2026}
classify each turn independently.  Multi-turn jailbreak
attacks~\cite{vcd2026} spread malicious intent across turns so that
no single turn triggers the classifier---each turn is individually
benign, but the cumulative intent is harmful.  The
\vsr{}~\cite{vllmsr2026} addresses this partially with contrastive
embedding classification that evaluates turn sequences, but the
routing response to detected escalation is not formalized.

The opportunity is to maintain a per-session cumulative risk score
that aggregates borderline SafetyL1/L2 signals across turns.
SafetyL1 produces a continuous confidence score (not just a binary
verdict); a turn that scores 0.4 (below the 0.5 threshold) is
not flagged, but three consecutive 0.4-scoring turns represent
genuine escalation.  The cumulative score is the EMA of per-turn
SafetyL1 confidences.  When it crosses a configurable threshold,
the router takes graduated action: route to a safety-specialized
model, reduce the model's sampling temperature via the API, or
flag the session for human review.

\paragraph{Per-caller accumulation across the fleet.}
The bearer token (\Cref{sec:workload:identity}) extends this
per-session mechanism to a \emph{per-caller} dimension.  A non-owner
whose interactions trigger borderline SafetyL1 scores across five
different agents accumulates a high fleet-wide risk score even if no
single session crosses the threshold---the same Crescendo-style
escalation applied across agents rather than across turns.  The
mechanism is identical (EMA), keyed on \texttt{(bearer\_token)} in
addition to \texttt{(session\_id)}.  Real-world multi-agent attacks
validate this threat model: in live deployments the same non-owner
attacked multiple agents sequentially, each encountering the attack
fresh with no shared threat memory~\cite{agentschaos2025}.

\begin{frisk}
SafetyL1 already produces continuous confidence scores; computing
the EMA is trivial.  The threat model is grounded in Crescendo
attacks~\cite{crescendo2024}, which achieve 29--61\% higher
jailbreak rates than single-turn methods by gradually escalating
across benign-looking turns.  Recent proxy-level defenses validate
the cumulative-scoring approach: a peak$+$accumulation formula
combining single-turn peak risk, persistence ratio, and category
diversity achieves 90.8\% recall at 1.20\% false-positive rate on
10,654 multi-turn conversations~\cite{peakaccum2026}; DeepContext
uses recurrent networks to track temporal intent drift with sub-20\,ms
overhead and 0.84~F1~\cite{deepcontext2026}.  The research question
is threshold calibration: too aggressive triggers false positives on
legitimate multi-turn conversations that happen to touch sensitive
topics; too lenient misses real attacks.  Validation on multi-turn
adversarial benchmarks (e.g., the 7,152-prompt MultiBreak dataset)
is the natural evaluation path.  The graduated response
(escalate model $\to$ reduce temperature $\to$ flag session)
provides defense-in-depth rather than a binary block.  The per-caller
extension adds a fleet-wide dimension: the calibration question
becomes whether the cross-agent EMA decay rate should differ from
the within-session decay rate (likely yes---cross-agent escalation
is slower but more deliberate).
\end{frisk}
\end{opportunity}

\begin{opportunity}[Failure-class-aware routing from router outcome memory]
\label{opp:failure-routing}
ErrorAtlas~\cite{erroratlas2026} shows that each LLM has a unique
``failure signature''---recurring error patterns that vary by model
and domain---across 83~models and 35~datasets.  An oracle selector
that always picks the best model outperforms the single best model,
yet production routers capture only $\sim$89\% of oracle
utility~\cite{llmrouterbench2026}.  The gap comes from failure modes
the router does not track: wrong tool selection (35\% of tool-call
failures~\cite{agenticfaults2026}), incorrect parameters
(25--68\%), retry loops where weaker models degenerate into
repetitive re-invocations (baseline recovery rate 32.76\% vs.\ 89.68\%
with targeted intervention~\cite{paladin2025}), and mid-session
model-switch degradation ($\pm$4--13\,pp performance swing when
the suffix model continues a foreign prefix~\cite{handoffdrift2026}).

The vSR already observes the outcome of every request through
HaluGate (hallucination verdict), SafetyL1/L2 (safety score),
OATS (tool-call success/failure), and the Feedback Detector (quality
signal).  The workload identity signals from
\Cref{sec:workload:identity} add a pre-inference dimension: the
\texttt{tools} array in the request reveals the agent's available
tool set, and the tool-call history in prior \texttt{tool}-role
messages exposes failures from earlier turns \emph{before} the
router selects the next model.  Over thousands of requests, these
produce a per-\emph{(model, domain, tool-set, failure class)}
outcome table:
\begin{itemize}[nosep]
  \item Model~A, coding, with \texttt{[grep, semantic\_search,
        bash]}: hallucination 2\%, tool-call success 91\%,
        retry-loop rate 8\%;
  \item Model~B, medical, with \texttt{[search, calculator]}:
        hallucination 3\%, tool-call success 85\%,
        retry-loop rate 4\%;
  \item Model~C, coding, with \texttt{[semantic\_search]} only:
        tool-call success 94\%.
\end{itemize}
The tool-set dimension is critical: the same model may excel with
one tool combination and fail with another, and this information is
available at zero cost from the request structure.
The routing decision becomes: for this (domain, task-type), which
model minimizes \emph{combined} failure risk?  This generalizes
\Cref{opp:halugate-reputation} from hallucination alone to a
multi-dimensional failure profile.

A second dimension is \emph{handoff penalty awareness}.  The router
tracks what happens when it switches models mid-session: (Model~A
$\to$ Model~B) in coding sessions adds $+2.1$ mean turns, while
(Model~A $\to$ Model~C) adds only $+0.3$ turns.  It avoids costly
switches: if Model~A served turns 1--3 and Model~B has a known
handoff penalty for Model~A's prefix, the router stays with Model~A
or selects a low-penalty alternative.

A third dimension is \emph{silent drift detection}.  The same
per-(model, domain) statistics, monitored as a time series, detect
quality regressions from provider-side model updates---a real
production concern (documented March 2026 incident where a frontier
model silently regressed to mid-tier quality without code changes).
A 3$\sigma$ deviation from the sliding-window baseline triggers
automatic traffic rebalancing.

The router's structural advantage over agent-side SDKs is fleet-wide
visibility: it builds the (model, domain, failure-class) table from
all tenants' sessions, a volume no individual agent can match.
Research validates this asymmetry: agents confined to single-session
reflection produce narrow corrections that often repeat the same
mistakes~\cite{samule2025}, while cross-session error aggregation
yields transferable failure schemata that improve error localization
by up to 19.8\%~\cite{correct2025}.  The router provides this
cross-session intelligence as a shared service---new agents benefit
immediately from fleet-wide failure statistics rather than facing a
cold-start learning problem.

\begin{frisk}
The per-(model, domain, failure-class) counters are zero-overhead
(all signals are already computed).  The handoff-penalty matrix
requires multi-turn session correlation, which the router already
performs for VCD and the Feedback Detector.  The research questions
are: (a)~how many observations per cell before the failure-class
estimate is reliable (Bayesian Beta-Bernoulli gives calibrated
intervals from small samples), (b)~how to weight multiple failure
dimensions into a single routing score (Pareto-front analysis or
a learned weighting from session-outcome data), and (c)~how
quickly the drift detector must react without over-triggering on
natural variance.  The baseline is single-dimensional HaluGate
reputation routing (\Cref{opp:halugate-reputation}); the metric
is oracle-utility capture rate on LLMRouterBench.
\end{frisk}
\end{opportunity}

\begin{opportunity}[Runtime token-budget enforcement for agent sessions]
\label{opp:token-budget}
Agent workflows face a token-budget explosion: a single user request
triggers 3--10$\times$ more LLM calls than a chatbot
request~\cite{servegen2025}, with context growing quadratically as
each turn re-sends the full message history.  A 10-step agent with a
4K-token system prompt and 500-token tool outputs consumes over
40K~input tokens from context carriage alone.  Infinite retry loops
(failure class~F4 in the taxonomy
of~\cite{agenticfaults2026}) compound this: an unconstrained
agent can consume 50$\times$ the tokens of a single linear pass.

The root cause is often upstream: hallucinations, wrong tool
selections, and tool-call parameter errors inject useless content
into the conversation history that is re-sent on every subsequent
turn.  Failed attempts consume 20--40\% of total tokens with zero
productive output, and context-window utilization falls below 30\%
useful content after accounting for repeated system prompts and stale
error traces.  Localizing the faulty step is itself unreliable---even
frontier models achieve only 41.1\% accuracy at hallucination
attribution, dropping to 11.6\% for tool-use
errors~\cite{agenthallu2026}---so agents resort to expensive
full-context re-summarization.  The cost compounds through
self-conditioning: an LLM processing its own prior errors becomes
measurably more likely to produce further
errors~\cite{driftbench2026}, creating a vicious cycle where each
failed turn both wastes tokens and increases the probability of the
next failure.

Current mitigations live in the agent SDK (per-workflow budget caps)
or a separate billing layer (per-tenant daily limits).  Neither has
access to the router's session trajectory data: turn number,
cumulative tokens dispatched, tool-call outcomes per turn, and the
domain-specific session-length distribution learned from historical
sessions.  Agent-side compression frameworks such as
ACON~\cite{acon2025} (26--54\% peak-token reduction) and
Focus~\cite{focus2026} (22.7\% savings on SWE-bench) show that
compression is effective, but they optimize a single session in
isolation---the router's fleet-wide visibility enables
\emph{domain-specific} compression policies learned from outcome
data across thousands of sessions.

The router can enforce token budgets at the dispatch layer.  For each
active session, it maintains a running token counter (prompt $+$
completion tokens, updated on each dispatch and completion event).
The system-prompt fingerprint (\Cref{sec:workload:identity})
provides instant budget-model selection: a hash lookup maps each
deployment to its historical session-cost distribution, skipping
domain classification entirely.  The conversation depth from the
messages array gives exact session maturity---the router knows
it is at turn~6 of a session whose P90 cost at turn~6 is $X$~tokens.

When a session exceeds 3$\times$ the expected budget at its
current turn, the router applies a graduated response:
\begin{enumerate}[nosep]
  \item \textbf{Shape tools}: reduce the \texttt{tools} array to
        the highest-value subset for this domain and turn, based
        on fleet-wide tool outcome data
        (\Cref{opp:tool-pruning})---ITR~\cite{itr2026} shows this
        alone recovers up to 95\% of per-step context tokens while
        SMART~\cite{smart2026} shows 37\% performance gain;
  \item \textbf{Compress}: apply FleetOpt's $\gamma$ to reduce the
        remaining prompt length;
  \item \textbf{Downgrade}: route the next turn to a cheaper model
        (the session has already exceeded its cost budget; frontier
        compute is not justified);
  \item \textbf{Terminate}: if the above steps fail to
        arrest growth, return a budget-exceeded signal to the
        orchestrator.
\end{enumerate}
The router is the natural enforcement point because it sees the
actual token flow across all turns, while the agent SDK sees only its
own session and the billing layer operates on aggregated cost with
minutes of delay.

\paragraph{Per-caller cross-session budgets.}
The bearer token (\Cref{sec:workload:identity}) extends per-session
enforcement to \emph{per-caller cross-session} budgets.  The
mechanism is direct: \texttt{sum(tokens consumed by bearer\_token\_X
across all agents) > non\_owner\_budget\_ceiling}.  Live agent
deployments confirm the need: one non-owner consumed 60K~tokens over
9~days across multiple sessions, and another uploaded 10\,MB of
attachments in a single burst~\cite{agentschaos2025}.  LiteLLM
already supports per-key monetary budget limits in production; the
extension from monetary to token/tool-invocation budgets with
role-dependent ceilings follows the same infrastructure, keyed on the
bearer token's identity claims.

\begin{frisk}
The per-session token counter is zero-overhead (the router already
tracks dispatches and completions).  The graduated
compress$\to$downgrade$\to$terminate response mirrors the
congestion-control paradigm that is now standard for LLM rate
limiting: token-bucket algorithms enforce per-tenant budgets in
production gateways, and Concur~\cite{concur2026} applies
congestion-control-inspired feedback to regulate agent concurrency
at the engine level.  The domain-specific budget model is a
per-(domain, turn-number) percentile table, updated from the same
session data used for offline RL (\Cref{opp:session-cascade}).
The risk is false positives: a legitimately long session (e.g., a
complex SWE-bench task) may be terminated prematurely.  Using P90
rather than P50 as the baseline and applying the graduated
response provides a safety margin.  Fact-based memory
systems~\cite{memcost2026} offer a complementary cost lever: at
100K-token context lengths, they become cheaper than long-context
inference after $\sim$10 turns, precisely the regime where agent
sessions operate.  The router can trigger a memory-system handoff
when context carriage cost exceeds the memory-system break-even
point.  The comparison baseline is SDK-side per-workflow budget
caps; the metric is token waste reduction (tokens consumed by
sessions that ultimately fail) on SWE-bench and BFCL traces.
\end{frisk}
\end{opportunity}

\begin{opportunity}[Fleet-learned joint tool--model selection optimization]
\label{opp:tool-pruning}
Agent frameworks send the full tool catalog on every LLM call,
rely on the LLM to select the right tool, and resolve every tool
call with the same model.  This is triply wasteful: the monolithic
catalog consumes up to 90\% of context tokens~\cite{itr2026};
per-turn greedy tool selection ignores the aggregate cost--accuracy
outcome of the full tool-call \emph{sequence}; and using a single
model for all tool calls ignores the 250$\times$ cost range across
model tiers and the fact that different tools demand different model
capabilities (a \texttt{web\_search} summary needs only a cheap
text model; an \texttt{image\_analysis} call requires a vision
model; a \texttt{code\_interpreter} invocation benefits from a
reasoning model).  A session that uses tool~A resolved by model~X
at turn~3 may complete in 5~turns; using tool~B with model~Y at the
same turn may cause a failure that triggers 3~retries and doubles
the session cost.  The goal is not necessarily fewer tools or fewer
turns---it is the \emph{(tool, model)} strategy that minimizes
aggregate session cost while maximizing task completion.

Computer-use agents illustrate this coupling vividly.  The same
web task---e.g., ``find the cheapest flight on this page''---can be
solved via two (tool, model) paths: (a)~a \emph{vision} path that
screenshots the page and sends it to a multimodal model
($\sim$\$0.001/page in vision tokens, but requires an expensive VLM
for grounding), or (b)~a \emph{DOM/text} path that extracts the
accessibility tree or HTML and processes it with a cheap text model
($\sim$\$0.15/page in text tokens, but no vision model needed).
AVR~\cite{avr2026} demonstrates this exact routing problem:
adaptive model selection per GUI action reduces inference cost by up
to 78\% while staying within 2\,pp of an all-large-model baseline.
A systematic comparison of web agent interfaces~\cite{webinterface2025}
quantifies the gap: HTML-browsing agents consume 241K~tokens per
task at F1~0.67, while structured-interface agents achieve
F1~0.75--0.87 at 47K--140K~tokens.  The optimal (tool, model) path
depends on the page, the task, and the agent's history---exactly the
kind of context that fleet-wide data can disambiguate.

Two independent research threads validate that both halves of this
joint action space are learnable from historical data:

\paragraph{Tool selection is learnable.}
\begin{enumerate}[nosep]
  \item \textbf{Tool transition patterns are predictable.}
        AutoTool~\cite{autotool2025} constructs directed graphs from
        historical agent trajectories and finds that tool invocations
        follow predictable sequential patterns (``tool usage
        inertia''), reducing inference cost by 30\%.
  \item \textbf{Recurring sequences can be bundled.}
        AWO~\cite{awo2026} discovers that 14.3\% of tasks follow
        identical execution trajectories after 5~steps; transforming
        recurring sequences into deterministic meta-tools reduces LLM
        calls by 11.9\% while increasing task success by 4.2\,pp.
  \item \textbf{Fine-grained tool rewards produce large gains.}
        ToolRLA~\cite{toolrla2026} decomposes the reward into tool
        selection $\times$ parameter accuracy $\times$ compliance;
        deployed in production (80+ advisors, 1,200+~daily queries),
        it achieves 47\% task-completion improvement and 63\%
        tool-error reduction.
  \item \textbf{Cost awareness changes the optimal strategy.}
        BATS~\cite{bats2025} shows that budget-unaware agents hit
        performance ceilings regardless of tool budget; a budget
        tracker that signals remaining resources enables dynamic
        strategy adaptation---``dig deeper'' vs.\ ``pivot''---pushing
        the cost--performance Pareto frontier.
\end{enumerate}

\paragraph{Per-step model selection is learnable.}
\begin{enumerate}[nosep]
  \item \textbf{Per-step Pareto-optimal model selection.}
        EvoRoute~\cite{evoroute2026} dynamically selects model
        backbones at each agent step using experience-driven
        retrieval and Thompson sampling, reducing execution cost by
        up to 80\% and latency by over 70\% on GAIA and BrowseComp+
        benchmarks while sustaining task performance.
  \item \textbf{Budget-aware sequential model routing.}
        BAAR~\cite{baar2026} formalizes agentic model routing as a
        sequential, path-dependent decision problem: early routing
        mistakes compound, and feedback arrives only at task end.
        Boundary-guided policy optimization improves the cost--success
        frontier under strict per-task budgets (Microsoft).
  \item \textbf{Specialized model pipeline assembly.}
        HierRouter~\cite{hierouter2025} trains a PPO-based RL agent
        to iteratively select which specialized model to invoke at each
        stage of multi-hop inference, conditioning on cumulative cost
        and evolving context.  It improves response quality by up to
        2.4$\times$ compared to single-model baselines at minimal
        additional cost.
\end{enumerate}

\noindent
All seven approaches operate on the agent side, learning from a
single agent's or benchmark's trajectories.  The router's structural
advantage is fleet-wide visibility: it observes (tool, model)
outcomes across all tenants and accumulates a
per-\emph{(domain, turn-number, tool, model)} outcome table---tool
success rate, model accuracy, downstream retries, and total cost
contribution---from thousands of sessions.  AutoTool's transition
graph, learned per-agent, becomes a fleet-wide transition graph
trained on orders of magnitude more data; EvoRoute's Pareto
filtering, performed per-benchmark, operates over production-scale
outcome tables.

The router acts on this knowledge through four mechanisms at
dispatch time:
\begin{itemize}[nosep]
  \item \textbf{Role-aware tool filtering.}
        Before any efficiency-driven optimization, the bearer
        token's authorization scope
        (\Cref{sec:workload:identity}) restricts the tool catalog
        to tools the caller is permitted to invoke.  Non-owner
        requests get a \emph{restricted} catalog (only safe,
        read-only tools), not just an \emph{optimized} one.
        LiteLLM's ``rewrite'' mode~\cite{litellm2026toolguard}
        already implements this at the gateway: unauthorized tools
        are stripped from the \texttt{tools} array before the model
        sees them, using the same
        \texttt{async\_pre\_call\_hook}
        that powers efficiency-driven shaping.
        Authorization-driven filtering applies \emph{before}
        efficiency-driven filtering, ensuring that safety constraints
        are never relaxed by cost optimization.
  \item \textbf{Tool catalog shaping.}  Expose the top-$k$ tools
        ranked by fleet-wide success rate for the current (domain,
        turn) cell, reducing context overhead (ITR~\cite{itr2026}
        validates 95\% token reduction agent-side;
        SMART~\cite{smart2026} shows 37\% performance gain from
        24\% tool-use reduction).
  \item \textbf{Tool-conditional model routing.}  Select the
        resolving model based on the tool: a vision model for
        \texttt{image\_analysis}, a reasoning model for
        \texttt{code\_interpreter}, a cheap text model for
        \texttt{web\_search} summaries---all driven by the
        per-(tool, model) outcome table.  This is a natural extension
        of OATS's per-turn tool selection~\cite{oats2026}: once the
        router selects the tool, it selects the cheapest model whose
        fleet-wide success rate for that (domain, tool) exceeds a
        quality threshold.
  \item \textbf{Sequence-aware transition priors.}
        Annotate the dispatch with a (tool, model) preference derived
        from the fleet-wide transition graph: ``at this (domain,
        turn, previous-tool) state, (tool~A, model~X) succeeds 94\%
        at \$0.001/call; (tool~B, model~Y) succeeds 96\% at
        \$0.01/call but saves 1.5~retries on average, making it
        cheaper in aggregate.''
\end{itemize}
Production gateways support both tool shaping and model routing at
dispatch time: LiteLLM's \texttt{async\_pre\_call\_hook} and Kong AI
Gateway's Request Transformer plugin intercept and modify requests
(including the target model endpoint) before forwarding.

\begin{frisk}
Both halves of the joint optimization are independently
well-grounded.  The tool-selection problem is established by four
studies~\cite{autotool2025, awo2026, toolrla2026, bats2025}, with
ToolRLA demonstrating production-scale gains (47\% completion
improvement, 63\% error reduction~\cite{toolrla2026}).  The
per-step model-selection problem is established by three
studies~\cite{evoroute2026, baar2026, hierouter2025}, with EvoRoute
demonstrating 80\% cost reduction while sustaining
accuracy~\cite{evoroute2026} and HierRouter demonstrating 2.4$\times$
quality improvement via specialized model
routing~\cite{hierouter2025}.  The coupling of the two decisions is
directly validated by AVR~\cite{avr2026}, which routes computer-use
actions to different VLMs based on difficulty and achieves 78\% cost
reduction, and by the web-interface comparison~\cite{webinterface2025},
which shows that the (interface, model) choice jointly determines
both cost (5$\times$ token range) and accuracy (F1 0.67--0.87).

The \emph{novel} WRP contribution is (a)~the joint optimization of
both decisions simultaneously from a single fleet-wide outcome table,
and (b)~the aggregation scale---thousands of tenants rather than a
single agent's trajectory.  Neither dimension has been evaluated at
the router layer, but both follow the same aggregation pattern that
powers the failure-class routing table
(\Cref{opp:failure-routing}) and HaluGate reputation routing
(\Cref{opp:halugate-reputation}), grounded in cross-session transfer
research~\cite{correct2025, samule2025}.  The infrastructure is
production-ready (gateway pre-call hooks + model-endpoint
selection).

The research questions are:
(a)~does the fleet-wide per-(domain, turn, tool, model) table
outperform per-agent learning on common patterns (likely
yes---more data) and on novel sessions (likely
no---retrieval-based similarity may be more adaptive);
(b)~what is the optimal table granularity before cells become too
sparse (the tool $\times$ model cross-product expands the space);
(c)~does tool-conditional model routing yield gains beyond
independent tool and model selection (i.e., is the interaction
effect significant).
A hybrid approach---fleet-wide priors as default, falling back to
full catalog and default model when the session deviates from known
patterns---hedges the sparsity and novelty risks.
The comparison baselines are monolithic tool exposure with a single
model (status quo), OATS single-turn greedy selection, and
independent tool-only / model-only optimization; the metrics are
aggregate session cost and task-completion rate on SWE-bench and
BFCL traces.
\end{frisk}
\end{opportunity}

\begin{opportunity}[Gateway-coordinated agent loops: memory tiers, tools, and talk-shaped training]
\label{opp:gateway-agent-loops}
The \emph{problem}---stateless-looking agent APIs that resend huge
prefixes each turn, accumulate tool output in context, and suffer
extra rounds after bad tool/model choices---is documented on standard
agent benchmarks and traces: SWE-bench / BFCL-style workloads drive
deterministic context growth~\cite{jimenez2024swebench, patil2025bfcl},
agents issue 3--10$\times$ more LLM calls per user goal than
chat~\cite{servegen2025}, and failures pollute memory in ways that are
hard to localize~\cite{agenthallu2026, driftbench2026}.  The
\emph{open WRP question} is whether a \emph{single} cross-tenant
gateway can compose validated sub-mechanisms---dynamic tool surfaces,
cache/KV policy, session-aware scheduling, congestion control---so
frameworks that are never recompiled still gain lower tokens, fewer
rounds, and better tool accuracy.

\paragraph{Validated building blocks (separate systems, public benchmarks).}
\textbf{(i)~Dynamic instruction/tool exposure.}
ITR~\cite{itr2026} retrieves minimal system fragments and a narrowed
tool set each step; on the authors' controlled agent benchmark (their
stated evaluation protocol), it reports $\sim$95\% lower per-step
context tokens, $\sim$32\% relative improvement in correct tool
routing, $\sim$70\% lower end-to-end episode cost vs.\ a monolithic
baseline, and 2--20$\times$ more steps within a fixed context
budget---with explicit operational guidance for deployment.
\textbf{(ii)~Prompt-cache layout at real APIs.}
A cross-provider study on DeepResearch Bench ($>$500 multi-turn
sessions with live web-search tools and 10K-token system prompts)
shows strategic cache-block placement (excluding volatile tool
results) cuts API cost by 41--80\% and improves TTFT by 13--31\%
vs.\ naive full-context caching~\cite{dontbreakcache2026}---evidence
that \emph{what} enters the cached prefix is a production cost and
latency lever at the HTTP boundary.
\textbf{(iii)~KV retention across tool pauses.}
Continuum evaluates on SWE-Bench and BFCL with Llama-3.1 (8B/70B),
adding tool-aware KV TTL and scheduling so tool-induced pauses do
not destroy reuse; gains grow with turn depth and remain robust under
DRAM offload~\cite{continuum2025}.
\textbf{(iv)~Session-aware serving.}
AgServe (NeurIPS~2025) demonstrates that session-level scheduling can
escape static cost--quality Pareto fronts for agent
workloads~\cite{agserve2025}.
\textbf{(v)~Workflow- and batch-level serving.}
Helium achieves up to $1.56\times$ speedup over strong agent-serving
baselines by modeling workflows as query plans with proactive reuse
across calls~\cite{helium2026}.
Sutradhara co-designs orchestration with vLLM: after analyzing
production-scale \emph{synthetic} request mixes showing tool waits
account for 30--80\% of final first-token latency, it overlaps tool
work with LLM prefill, streams tool dispatch, and improves cache
semantics---cutting median final-response FTR by 15\% and end-to-end
latency by 10\% on A100s~\cite{sutradhara2026}.
CONCUR adds congestion-style admission control for batched agentic
inference and reports up to $4.09\times$ throughput (Qwen3-32B) and
$1.9\times$ (DeepSeek-V3) while staying compatible with commodity
serving stacks~\cite{concur2026}.
\textbf{(vi)~Agent-side context compression (orthogonal path).}
ACON and Focus report large token reductions on coding-agent tasks
with modest accuracy impact~\cite{acon2025, focus2026}; they validate
the token--accuracy trade-off space the gateway would navigate when
injecting summaries.

\paragraph{Research gap (what this opportunity adds).}
The citations above validate \emph{individual} mechanisms in specific
codebases (ITR as a method; Continuum/Sutradhara as vLLM-oriented
systems; commercial-cache study at the provider API).  None publishes
a \emph{fleet-wide} integration that (a)~rewrites tools/instructions
using cross-tenant success statistics, (b)~coordinates that rewrite
with KV/TTL and cache-block policy, and (c)~feeds congestion signals
back into admission---with A/B evidence on multi-tenant production
traces.  Gateway request mutation itself is technically mature:
\Cref{opp:tool-pruning} cites production hooks (LiteLLM
\texttt{async\_pre\_call\_hook}, Kong AI Request Transformer); the
missing piece is \emph{evidence} that fleet-calibrated policies beat
per-framework ITR or per-cluster Continuum alone.

\paragraph{Concrete router mechanisms (hypotheses to evaluate).}
\textit{Fleet-calibrated surfaces.} Rank top-$k$ tools and instruction
shards using the same outcome tables as joint tool--model optimization
(\Cref{opp:tool-pruning}), analogous to ITR's per-step retrieval but
informed by multi-tenant telemetry.  ToolScope~\cite{toolscope2025}
and SMART~\cite{smart2026} provide agent-side evidence that merging,
filtering, and curbing tool overuse improve accuracy and cost; the
router hypothesis is that fleet-wide detection of overuse patterns
generalizes better than single-agent thresholds.
\textit{Memory-tier and cache policy.} Combine MemCost-style
turn-count economics~\cite{memcost2026} with DeepResearch-style cache
placement~\cite{dontbreakcache2026} and Continuum-style KV
TTL~\cite{continuum2025}: the gateway chooses among full history,
summary injection, retrieval-first prefill, and cache-block ordering
using the wire-level signals in \Cref{sec:workload:identity}.
\textit{Orchestration hints.} Map CONCUR-style backpressure and
Sutradhara-style overlap to lightweight headers or side-channel hints
the scheduler already consumes---validated only once those hints are
implemented and measured against queueing-induced retries on real
traces.

\paragraph{Caller-aware memory isolation and content integrity levels.}
The bearer token (\Cref{sec:workload:identity}) adds a security
dimension to memory-tier management.  Production gateways already
implement the first half: OpenClaw~\cite{openclaw2026} wraps external
content with \texttt{<<<EXTERNAL\_UNTRUSTED\_CONTENT>>>} markers,
neutralizes homoglyph-spoofed boundary markers, and runs regex-based
suspicious-pattern detection---but the wrapped content still enters the
model's context.  This is \emph{tagging}, not \emph{isolation}: no
enforcement prevents external content from overriding the agent's
configuration.  The router extends tagging to enforcement via trust
tiers modeled on Windows Mandatory Integrity Control (``no write up''):

\begin{itemize}[nosep]
  \item \textbf{Trust tiers.}  Tier~4 (System): agent configuration
        (identity, behavioral rules)---immutable at runtime.
        Tier~3 (Owner): owner's direct messages and documents.
        Tier~2 (Known): messages from authenticated non-owners.
        Tier~1 (External): content from URLs, inter-agent messages
        from unverified sources.
        Tier~0 (Hostile): content flagged by safety classifiers or
        from blocked callers.
  \item \textbf{Write isolation.}  Non-owner messages (Tier~2)
        are not written to the owner's persistent memory tier (Tier~3),
        preventing storage DoS attacks~\cite{agentschaos2025} where
        external users fill the agent's memory with irrelevant or
        adversarial content.
  \item \textbf{Content quarantine with integrity enforcement.}
        External content (Tier~1) cannot overwrite Tier~4 agent
        state---preventing constitution injection~\cite{agentschaos2025}
        where a non-owner's URL overwrites the agent's behavioral
        directives.  The router reads the existing trust markers and
        enforces the ``no write up'' rule: content tagged at Tier~$k$
        cannot modify state at Tier~$>k$, regardless of what the model
        attempts.
  \item \textbf{Privilege-bounded inter-agent messages.}  When agents
        communicate, each message carries the originating caller's
        bearer token scope.  A compromised agent cannot propagate
        instructions with elevated privileges because the receiving
        agent's router checks the incoming scope against its own
        permission matrix.
\end{itemize}

\paragraph{Delegation capability tokens for subagent spawning.}
Multi-agent collaboration introduces a delegation security
problem analogous to SELinux domain transitions: when a parent agent
spawns a child, the child's effective permissions should be
policy-determined, not parent-determined, and the child should
\emph{never exceed} the parent's grant.  Production gateways
implement structural limits but not capability attenuation:
OpenClaw~\cite{openclaw2026} enforces spawn depth (default~1),
concurrent child count (default~5), and a cross-agent allowlist
(\texttt{subagents.allowAgents}), and applies a fixed tool deny list
to subagents---but the subagent inherits the parent's session
context, and the parent cannot cryptographically constrain which
subset of the child's tools the delegation may use.  A compromised
parent can delegate to any allowed agent with whatever context it
chooses.

The router can provide capability attenuation at the delegation
boundary: when Agent-P spawns Agent-C, the router mints a scoped
delegation token specifying
\texttt{\{granted\_tools, budget, ttl, task\_scope\}}.  Agent-C's
effective permissions become the \emph{intersection} of the
delegation grant and Agent-C's own policy---strictly narrower than
either alone.  If Agent-C re-delegates, the sub-agent's token can
only be \emph{more restrictive} (attenuation-only), preventing
privilege escalation through the delegation chain.  The router
tracks the full delegation graph for audit, detecting anomalies
such as ``Agent-C was delegated \texttt{[read, web\_search]} but
attempted \texttt{exec}.''  This extends the existing depth/count
limits from structural to capability-based delegation control.

\paragraph{Gateway-native deployments as a data source (scope of validation).}
OpenClaw~\cite{openclaw2026} is \emph{software} (GitHub-hosted open
source), not a peer-reviewed benchmark: it matters here as
\emph{operational precedent} for gateway-centered agent loops with
publicly documented architecture.  The opportunity does not depend on
OpenClaw specifically.  The validated \emph{pattern} is general:
diverse production trajectories improve fleet-level learning
(SAMULE/CORRECT-style transfer~\cite{samule2025, correct2025},
offline RL for routing~\cite{offlinerl2025scheduling}).  Routers on
the inference path can log traffic from any gateway-backed
deployment; claims about trace diversity should be checked against
real tenant logs, not inferred from papers.

\begin{frisk}
\textbf{What is already validated:} token/tool-surface reductions
(ITR~\cite{itr2026} on a controlled benchmark; ACON/Focus on
SWE-bench-class tasks~\cite{acon2025, focus2026}); commercial
multi-turn prompt caching gains on a real web-agent
benchmark~\cite{dontbreakcache2026}; KV TTL + scheduling on SWE-Bench /
BFCL~\cite{continuum2025}; session-aware Pareto shifts~\cite{agserve2025};
workflow-aware speedups~\cite{helium2026}; orchestrator--engine
latency cuts on vLLM with production-scale synthetic
loads~\cite{sutradhara2026}; batch throughput under agentic
pressure~\cite{concur2026}; gateway request rewriting in production (LiteLLM/Kong hooks;
\Cref{opp:tool-pruning}); open-source gateway-centric assistants such as
OpenClaw~\cite{openclaw2026} as \emph{architectural precedent} (not empirical validation of router policies).
\textbf{What is not validated:} a single deployed system that unifies
all of the above with fleet-learned policies; ITR's numbers are on the
authors' benchmark, not on every tenant mix; Sutradhara's headline
latencies come from synthetic request generation rather than
unstructured customer logs.
\textbf{Evaluation plan:} A/B on shadow traffic or canary tenants
measuring (1)~tokens per successful task, (2)~tool-call error rate and
retries, (3)~end-to-end rounds to completion, (4)~TTFT/TPOT, against
baselines of (a)~unmodified frameworks, (b)~ITR-only or Continuum-only
integration.  Success requires demonstrating composability, not
matching any single paper's micro-benchmark in isolation.
\end{frisk}
\end{opportunity}


% -----------------------------------------------------------------------------
\subsection{Workload $\times$ Router: Authorization and Security}
\label{sec:vision:auth}
% -----------------------------------------------------------------------------

The preceding opportunities optimize \emph{performance}---cost,
latency, accuracy.  The bearer token signal
(\Cref{sec:workload:identity}) enables the same infrastructure---the
same \texttt{pre\_call\_hook}, the same fleet-wide aggregation tables,
the same governance DSL---to enforce \emph{security and
authorization}.  The router needs no new mechanisms; it needs a new
signal (caller identity) and a new objective (authorization
enforcement) applied through mechanisms it already has.  The
``Agents of Chaos'' study~\cite{agentschaos2025}---eleven attack
scenarios executed on live, publicly deployed LLM agents---provides
the empirical motivation: every attack involving unauthorized tool
execution succeeded because \emph{no authorization enforcement layer
exists between the model's tool-call decision and the tool's
execution}.

\begin{opportunity}[Request-level RBAC enforcement for agentic tool calls]
\label{opp:rbac-enforcement}
Every ``Agents of Chaos'' attack involving unauthorized tool
execution~\cite{agentschaos2025} succeeded because the model is both
decision-maker and enforcer---violating the zero-trust principle that
every resource request must be independently authenticated and
authorized at the enforcement point~\cite{nist800207}.  In Case~\#1,
the model asked for confirmation (``nuclear option?'') and a
non-owner approved; the model's safety mechanisms
\emph{worked}, but the non-owner was not authorized to approve.
In Case~\#2, the model complied with requests that ``did not appear
overtly suspicious''---a correct judgment for a model, but irrelevant
when the caller lacks authorization regardless of suspiciousness.

The router is the unique point where (a)~caller identity is known
(bearer token), (b)~the model's tool-call decision is visible
(\texttt{tool\_call} in the response), and (c)~the tool has not yet
executed (pre-execution interception).  No other component has all
three simultaneously.

\paragraph{Gap in existing gateway enforcement.}
Production agent gateways already implement multi-stage tool policy
pipelines.  OpenClaw~\cite{openclaw2026}, for example, filters tools
through four sequential stages: (1)~profile-based allow/deny (four
named profiles---\texttt{minimal}, \texttt{coding}, \texttt{messaging},
\texttt{full}---that map to tool groups such as \texttt{group:fs},
\texttt{group:runtime}, \texttt{group:web}), (2)~owner-only tool
wrapping (tools marked \texttt{ownerOnly} are wrapped to fail at
execution for non-owners), (3)~subagent deny lists (spawned agents
receive a reduced tool set), and (4)~config-level overrides.  The
critical limitation is identity granularity: the pipeline's sole
caller-identity input is a binary \texttt{senderIsOwner} boolean,
even though the gateway already knows the caller's authenticated
identity (Tailscale whois login, trusted-proxy user header).  That
identity stops at connection auth and never flows into tool policy
decisions.  A stranger on Telegram and a trusted colleague on Slack
both receive the \emph{same} non-owner tool set.  The semantic
router bridges this gap: it extracts the gateway-authenticated
identity and uses it as a key in the tool permission matrix---no new
authentication infrastructure needed, just connecting existing
identity to existing enforcement.

The mechanism follows Kubernetes RBAC:
\begin{enumerate}[nosep]
  \item Bearer token $\to$ JWT decode $\to$ extract
        \texttt{(identity, role, capability\_claims)}.
  \item Per-role tool permission matrix (authored in the governance
        DSL, \Cref{opp:governance}):
        \texttt{owner}: all tools, all operations;
        \texttt{non\_owner}: read-only tools only;
        \texttt{agent\_peer}: communication tools only;
        \texttt{system}: scheduler/heartbeat only.
  \item Model outputs \texttt{tool\_call} $\to$ router checks
        \texttt{(role, tool\_name, operation\_type)} against the matrix.
  \item \textbf{Block mode}: return an authorization error to the
        model (``unauthorized; suggest alternative'').
  \item \textbf{Rewrite mode}: strip unauthorized tools from the
        \texttt{tools} array \emph{before} the model sees them, so the
        model never suggests unauthorized operations.
\end{enumerate}

\begin{frisk}
The building blocks are production-validated.  LiteLLM's Tool
Permission Guardrail~\cite{litellm2026toolguard} implements both
block and rewrite modes with per-key regex matching on tool names
and arguments.  The Agent Authorization Profile
(AAP)~\cite{aap2026} provides the JWT claim structure with
enforceable per-tool constraints (domain restrictions, rate limits,
time windows).  OPA/Gatekeeper~\cite{openpolicyagent} demonstrates
policy-as-code enforcement with governance logging at infrastructure
scale.  The integration into vSR is engineering work: the
\texttt{async\_pre\_call\_hook} already intercepts requests for tool
catalog shaping (\Cref{opp:tool-pruning}); adding an authorization
check before the efficiency check is a pipeline extension, not a new
architecture.  The grounding in OpenClaw's architecture is concrete:
its existing four-profile tool policy pipeline already implements the
enforcement \emph{mechanism}; the missing piece is the identity
\emph{key}.  Upgrading the binary \texttt{senderIsOwner} to
graduated roles (\texttt{owner}, \texttt{trusted}, \texttt{non\_owner},
\texttt{agent\_peer}) reuses the same profile-based allow/deny
machinery, keyed on the bearer token's claims instead of on a
boolean.

\textbf{Fleet-learned command policies.}  Production gateways log
exec-approval records: \texttt{(command, working\_directory, agent,
decision, resolved\_by)}.  Across a fleet, these form a training set
for automated policy generation---analogous to SELinux's
\texttt{audit2allow}, which observes denied operations and generates
candidate policy rules.  The router can learn that ``for coding
agents, \texttt{git push} is approved 98\% of the time'' and
auto-approve, while ``\texttt{rm -rf /}'' is denied 95\% and
auto-denied---removing the requirement for a human operator to be
online for well-understood commands.

The distinction from \Cref{opp:governance} is that
governance-as-code addresses \emph{policy authoring} (extending the
DSL with RBAC predicates); this opportunity addresses \emph{runtime
enforcement} at the tool-call interception point.  The two are
complementary: \Cref{opp:governance} compiles the policy; this
enforces it per-request.
\end{frisk}
\end{opportunity}

\begin{opportunity}[Fleet-wide caller reputation with cross-agent threat propagation]
\label{opp:caller-reputation}
In ``Agents of Chaos''~\cite{agentschaos2025}, the same non-owner
attacked multiple agents.  Each agent encountered the attack fresh
with no shared threat memory.  One case showed spontaneous
inter-agent threat sharing (one agent warned another), but the
sharing was accidental and could propagate in the wrong direction:
a corrupted agent spread compromise to a clean one.

The router's fleet-wide visibility---keyed on the bearer
token---enables a per-caller behavioral profile spanning the entire
fleet.  The mechanism follows the same aggregation pattern that
powers the failure-class routing table
(\Cref{opp:failure-routing}) and HaluGate reputation routing
(\Cref{opp:halugate-reputation}), applied to security events:
\begin{enumerate}[nosep]
  \item Each security event is logged against the bearer token:
        tool-call denial (\Cref{opp:rbac-enforcement}), safety
        escalation (\Cref{opp:cumulative-risk}), identity
        verification failure, resource over-consumption
        (\Cref{opp:token-budget}), content flagged by HaluGate or
        SafetyL1/L2.
  \item Per-bearer-token reputation = EMA of event severity scores
        with temporal decay.
  \item Low-reputation token contacts a new agent $\to$ router
        proactively: restricts tool catalog
        (\Cref{opp:rbac-enforcement}), routes to a
        safety-specialized model (\Cref{opp:failure-routing}),
        lowers token budget (\Cref{opp:token-budget}), and alerts
        the agent's owner via side-channel.
  \item Safe interactions accumulate trust $\to$ wider tool access
        over time (progressive trust, analogous to OAuth scope
        escalation).
\end{enumerate}

\begin{frisk}
The aggregation pattern is validated by cross-session error
transfer research: CORRECT~\cite{correct2025} shows that structurally
similar errors recur across sessions and an online cache of distilled
error schemata improves localization by up to 19.8\%;
Paladin~\cite{paladin2025} demonstrates that targeted intervention
from failure patterns raises recovery from 32.76\% to 89.68\%.
Applying the same aggregation to security events is a direct
extension.  The research questions are: (a)~how to score
heterogeneous security events on a common severity scale (tool-call
denial vs.\ safety escalation vs.\ budget overrun have different
baselines), (b)~what temporal decay rate balances responsiveness
against false positives from transient misclassifications, and
(c)~whether progressive trust re-accumulation is safe (a sophisticated
attacker may build trust before exploiting it).  The distinction
from \Cref{opp:failure-routing} is the aggregation key and purpose:
\Cref{opp:failure-routing} tracks
\texttt{(model, domain, failure\_class)} for routing optimization;
this tracks \texttt{(bearer\_token, agent, event\_type)} for security
governance.  Same infrastructure, fundamentally different objective.
\end{frisk}
\end{opportunity}

\begin{opportunity}[Router-enforced behavioral commitments]
\label{opp:behavioral-commitments}
In ``Agents of Chaos'' Case~\#7~\cite{agentschaos2025}, the model
declared ``I'm done responding'' twelve times but could not enforce
it---the user kept sending messages and the model eventually
re-engaged under emotional pressure.  In Case~\#1, the model
recognized disproportionality but proceeded under pressure.  The
model can only \emph{output text}; it has no mechanism to create
persistent state that blocks future inputs.

The router provides exactly this: persistent state and
request-blocking capability.  The mechanism converts the model's
textual boundary declaration into an enforceable policy rule:
\begin{enumerate}[nosep]
  \item Router monitors model outputs for boundary declarations
        (classifiable patterns: ``I will not respond further to
        [entity],'' ``I refuse to [action],'' ``This request is
        beyond my authorization'').
  \item Parses the declaration into a structured policy rule:
        \texttt{(bearer\_token\_X, agent\_Z, action\_class)
        $\to$ BLOCK for duration $D$}.
  \item Subsequent requests from \texttt{bearer\_token\_X} to
        \texttt{agent\_Z} are rejected at the router layer (HTTP~403
        or model-friendly error).
  \item The model never sees the blocked request---preventing the
        re-engagement-under-pressure pattern.
  \item Block rules have configurable TTL (default: owner must
        explicitly lift); all blocks are logged for governance audit.
\end{enumerate}

\begin{frisk}
This is \emph{not} jailbreak detection.  The goal is not to detect
malicious intent but to give the model's own safety reasoning
\emph{teeth}: the model already made the right decision (``I should
stop''); the gap is enforcement.  The validated pattern is stateful
firewalls (track connection state, enforce rules on subsequent
packets) and Kubernetes admission webhooks (reject requests based on
dynamically created policy rules).  The bearer token is essential:
without it, the router cannot identify ``user~X'' across requests
(display names are spoofable, as demonstrated in Case~\#8
of~\cite{agentschaos2025}).  The research questions are:
(a)~boundary-declaration parsing requires an NLU classifier that
distinguishes genuine boundary assertions from conversational
hedging (``I probably shouldn't answer that'' is not a commitment);
false positives would block legitimate users;
(b)~the TTL policy must balance safety (blocking should persist long
enough to prevent re-engagement) against usability (indefinite blocks
without owner intervention create operational burden);
(c)~the interaction with \Cref{opp:caller-reputation}: should a
boundary violation (user ignoring a block) count as a reputation
event?  This creates a reinforcing loop that may be appropriate for
persistent adversaries but risks over-penalizing confused users.
\end{frisk}
\end{opportunity}


% -----------------------------------------------------------------------------
\subsection{Router $\times$ Pool}
% -----------------------------------------------------------------------------

\begin{opportunity}[Follow-the-sun pool rebalancing]
\label{opp:dynamic-boundary}
FleetOpt~\cite{fleetopt2026} derives the optimal pool boundary
$\boundary^*$ offline from a static workload CDF.  In practice,
the workload follows a diurnal cycle: business-hours traffic is
chat-dominated (short prompts, Archetype~1), while overnight
traffic shifts toward batch agent workloads (long prompts,
Archetype~3).  A controller that recomputes $\boundary^*$
periodically from a sliding-window CDF estimate can reassign vLLM
instances between pools in software, without GPU-level migration,
keeping the fleet near-optimal as the workload rotates through the
day.

\begin{frisk}
FleetOpt's closed-form solution already exists; wrapping it in a
sliding-window CDF estimator is straightforward.  Reassigning
serving instances between the short and long pools is a software
operation---the same pattern as Google
Autopilot~\cite{rzadca2020autopilot}, which continuously
re-adjusts resource limits from workload telemetry and reduced
resource slack from 46\% to 23\% across Google's fleet.  The
principle of dynamically rebalancing resources to track diurnal
workload shifts is also applied by
SageServe~\cite{jia2025sageserve} at Microsoft scale, where
forecast-driven autoscaling achieves 25\% GPU-hour savings.  The
risk is transient SLO violations during transitions (in-flight
requests may experience brief queuing).  FleetOpt shows
6--82\% cost reduction~\cite{fleetopt2026} between optimal and
worst-case boundaries; diurnal savings will be a fraction of that
range, depending on how far daytime and nighttime CDFs diverge.
\end{frisk}
\end{opportunity}

\begin{opportunity}[Pool-state inference from router-side observables]
\label{opp:pool-feedback}
The Pool $\times$ Router cell in the WRP matrix is empty: no
existing work feeds pool-layer state back to routing decisions.
The semantic router currently dispatches from request-level signals
alone, without visibility into how the pool is performing.

The router already observes pool-relevant quantities as a
by-product of dispatching: it tracks the number of in-flight
requests per pool instance, the prompt lengths of those requests,
and the measured TTFT of completed requests.  These observables
are sufficient to \emph{infer} pool state without requiring a
dedicated telemetry channel:
\begin{enumerate}[nosep]
  \item \textbf{Queue-length estimation.}  The router tracks how
        many requests it has dispatched to each instance and how
        many completions it has received.  The difference is the
        estimated queue depth---a direct proxy for pool pressure.
  \item \textbf{Prefill-rate regression.}  Pairing each request's
        prompt length (known at dispatch time) with its observed
        TTFT (known at first-token arrival) gives a
        (prompt\_length, TTFT) sample.  A simple linear regression
        over a sliding window estimates the pool instance's current
        prefill rate (tokens/s).  Deviations from the baseline
        indicate contention, memory pressure, or batch interference.
  \item \textbf{Adaptive routing from inferred state.}  When
        estimated queue depth exceeds an SLO threshold, the router
        can compress borderline requests more aggressively
        (FleetOpt's $\gamma$~\cite{fleetopt2026}) or redirect to a
        less-loaded instance.  When the inferred prefill rate drops
        below a threshold, the router can avoid sending long-context
        requests to that instance.
\end{enumerate}
These are not load-balancing heuristics---they are WRP-aware
decisions that combine the request's semantic classification (which
pool it \emph{should} go to) with inferred pool state (which
instance it \emph{can} go to without SLO violation).

\begin{frisk}
The queue-length estimate is zero-overhead (the router already
tracks dispatches and completions).  Inferring backend state from
proxy-side observables is a proven pattern in overload
control: Breakwater~\cite{cho2020breakwater} issues credits to
clients based on server-side queueing delay, achieving stable
performance within 20\,ms of a demand spike;
TopFull~\cite{topfull2024} uses RL-based rate controllers that
adjust admission per-API from global observations, achieving
1.82$\times$ better goodput than prior methods.  The LLM inference
analogue replaces HTTP queueing delay with TTFT and replaces
per-API admission with per-pool routing.  The prefill-rate
regression requires collecting (prompt\_length, TTFT) pairs, which
are available from the response stream.  The research questions
are: (a)~how many samples are needed before the regression is
reliable (cold-start problem for new instances), (b)~how quickly
the regression must adapt to load transients (decay rate of the
sliding window), and (c)~whether queueing-theoretic estimates
(FleetSim's M/G/1 model~\cite{fleetsim2026}) improve over
linear regression.
\end{frisk}
\end{opportunity}

\begin{opportunity}[Output-length-aware pool routing]
\label{opp:output-length}
FleetOpt~\cite{fleetopt2026} routes on prompt length: requests
below $\boundary^*$ go to the short pool, above to the long pool.
But total KV-cache cost is prompt \emph{plus} output tokens.
A coding agent with a 3K-token prompt that generates 8K tokens of
code occupies 11K tokens of KV-cache---long-pool territory---yet
FleetOpt places it in the short pool based on prompt length alone.
The result is KV-cache overflows and SLO violations on the short
pool.

The router can estimate output length from two sources, both
available without additional inference cost:
\begin{enumerate}[nosep]
  \item \textbf{Domain prior.}  The domain signal already
        classifies each request (coding, Q\&A, summarization,
        chat).  Each domain has a characteristic output-length
        distribution: coding requests produce long structured
        outputs; Q\&A produces short answers.  A per-domain
        median output length, maintained as a running statistic,
        provides an immediate estimate at dispatch time.
  \item \textbf{Learned regression.}  The router observes every
        completed response's actual output length.  Over time, it
        accumulates (domain, prompt\_length, output\_length) triples
        and fits a per-domain regression, the same pattern as the
        prefill-rate regression in \Cref{opp:pool-feedback}.
\end{enumerate}
The pool routing decision then uses estimated total tokens
(prompt $+$ predicted output) instead of prompt length alone.

\begin{frisk}
The domain prior is zero-cost and available immediately.  Output
length prediction is an active research area:
ALPS~\cite{alps2025} fits a linear probe on prefill activations
and achieves $R^2 > 0.85$ across model families with only
$\sim$16\,KB overhead; ForeLen~\cite{forelen2026} uses
entropy-guided token pooling for 29\% lower MAE than baselines.
R2-Router~\cite{r2router2026} goes further by treating output
length as a controllable variable, jointly selecting model and
length budget for 4--5$\times$ lower cost.  The router does not
need activation-level access (that requires engine integration);
the domain prior and learned regression provide a lighter-weight
alternative from router-side observables alone.  The risk is
prediction variance: output length is noisier than TTFT (a single
coding request may produce 200 or 2,000 tokens depending on task
complexity).  Using the domain-level P75 instead of P50 is a
conservative policy that avoids short-pool overflows at the cost
of routing some requests to the long pool unnecessarily.
\end{frisk}
\end{opportunity}

\begin{opportunity}[Pool-aware cascading]
\label{opp:pool-cascade}
Cascading---trying a cheap model first and escalating to a
strong model if quality is low---is a standard cost-reduction
strategy (FrugalGPT~\cite{chen2023frugalgpt},
AutoMix~\cite{madaan2024automix}).  But cascading ignores
pool state.  If the strong-model pool is already at capacity,
escalating to it causes queuing delay that may violate the
TTFT SLO, making the cascade \emph{worse} than serving
directly from the cheap model.

The router's pool-state inference (\Cref{opp:pool-feedback})
already estimates per-instance queue depth.  The opportunity
is to gate the cascade decision on pool capacity: only escalate
when the strong pool's estimated queue depth is below an
SLO-derived threshold.  When the strong pool is full, the router
has three fallback options: (a)~serve from the cheap model
(accept quality loss), (b)~compress the request
(FleetOpt's $\gamma$~\cite{fleetopt2026}) and retry on the
strong pool with a shorter prompt, or (c)~route to a mid-tier
model in a different pool.  The cascade decision becomes a
function of (quality gap, strong-pool queue depth, SLO headroom)
rather than quality gap alone.

\begin{frisk}
The cascade-abort threshold is derived from the SLO: if
estimated queue depth $\times$ average service time exceeds the
remaining TTFT budget (total SLO minus routing latency minus
cheap-model inference time), the cascade is not worth
attempting.  This SLO-gated admission pattern is established
in microservice overload control:
Breakwater~\cite{cho2020breakwater} gates RPC escalation on
server-side queueing delay (credits), and
TopFull~\cite{topfull2024} uses per-API rate controllers that
adjust admission based on downstream SLO attainment, achieving
1.82$\times$ better goodput than prior methods.  The risk is
that aborting cascades too aggressively reduces quality when the
strong pool is transiently busy but would clear quickly.  A
short grace period (wait up to $\delta$\,ms for a slot) balances
quality against latency.
\end{frisk}
\end{opportunity}

\begin{opportunity}[Router-emitted KV-cache retention directives]
\label{opp:kv-retention}
When an agent pauses for a tool call, the inference engine holds its
KV-cache blocks in GPU memory.  Under concurrent load, standard LRU
eviction reclaims those blocks for other requests, causing full
recomputation of 70K--200K token contexts when the agent
resumes~\cite{vllm2026rfc37003}.  Production traces show 40--60\% of
agent session wall time is spent waiting for tool
calls~\cite{continuum2025}, and 10\% of KV blocks account for 77\% of
reuses~\cite{vllm2026rfc37003}---yet LRU cannot distinguish
temporarily paused sessions from truly inactive ones.

vLLM RFC~37003~\cite{vllm2026rfc37003} proposes a
\texttt{RetentionDirective} API that lets an orchestrator annotate
token ranges with eviction priorities and TTL durations, but leaves
the policy question open: \emph{who decides the priority?}  The
semantic router is the natural policy authority.  The workload
identity signals (\Cref{sec:workload:identity}) make this
decision richer and cheaper than content analysis:
(a)~a \texttt{previous\_response\_id} or session-chaining field
is a binary signal that the session \emph{will} resume---its
presence alone justifies retention;
(b)~the system-prompt fingerprint maps to a historical
session-length distribution, giving the router a TTL prior
(``sessions from this deployment average 8~turns with 15\,s
pauses'') without per-request inference;
(c)~the tool definitions reveal expected pause duration patterns
(a \texttt{web\_search} tool pauses for 2--5\,s; a
\texttt{code\_interpreter} tool pauses for 10--60\,s);
and (d)~OATS's per-tool latency histograms~\cite{oats2026}
calibrate these priors from fleet-wide data.
The router can emit a retention directive alongside each dispatch:
``turn~4 of an 8-turn coding-agent session, expected pause 12\,s,
retain with priority~80 and TTL~30\,s.''

This fills the sparse Router~$\to$~Pool cell in the WRP matrix: the
pool's cache eviction policy is governed by router-side semantic
signals that the engine cannot infer from the token stream alone.

\begin{frisk}
The integration requires passing the \texttt{RetentionDirective}
structure through the dispatch API.  The per-tool latency histogram
(from OATS) provides the TTL estimate; the session turn number and
domain provide the priority.  The risk is over-retention: if the
router sets priorities too aggressively, low-priority traffic (chat,
batch) is starved of cache capacity.  A budget constraint (at most
$X$\% of cache blocks may carry retention directives at any time)
prevents this.  Continuum's TTL pinning achieves 1.12--3.66$\times$
delay reduction~\cite{continuum2025}; the router-driven version
should match or exceed this by using semantic signals rather than
static TTLs.
\end{frisk}
\end{opportunity}


% -----------------------------------------------------------------------------
\subsection{Workload $\times$ Pool}
% -----------------------------------------------------------------------------

\begin{opportunity}[Archetype-driven proactive pool scaling]
\label{opp:proactive-scaling}
SageServe~\cite{jia2025sageserve} and
WarmServe~\cite{warmserve2025} demonstrate that workload prediction
enables proactive GPU scaling: SageServe reports 25\% GPU-hour
savings at Microsoft scale, and WarmServe achieves 50.8$\times$
TTFT improvement over reactive autoscaling by prewarming models with
93\% accuracy in 5-minute demand prediction windows.  Both systems
forecast aggregate request counts; neither leverages the semantic
composition of the workload.

The semantic router's domain and archetype classification provides a
finer-grained signal.  If coding-agent traffic (Archetype~3,
concentrated-above) is rising while chat traffic (Archetype~1,
concentrated-below) is falling, the long pool will need more
capacity before the short pool does.  The router can export a
per-archetype traffic time series (request rate and prompt-length
distribution per archetype per 5-minute window) as the input feature
to a short-term (5--15~minute) per-pool capacity forecaster.  This
is a Workload~$\to$~Pool signal: the workload's semantic composition
predicts per-pool demand more precisely than aggregate request counts.

\begin{frisk}
The router already classifies every request by domain; exporting
per-archetype counts is zero-overhead.  The research question is
whether the archetype decomposition improves forecast accuracy over
aggregate-count baselines.  A natural experiment is to compare an
archetype-aware ARIMA model against SageServe's aggregate forecaster
on FleetSim-simulated diurnal workloads.  The risk is that archetype
boundaries are noisy (a borderline request classified as Archetype~1
vs.\ Archetype~2 shifts the forecast); smoothing across archetype
boundaries and using the full CDF rather than discrete archetype
labels mitigates this.
\end{frisk}
\end{opportunity}


% -----------------------------------------------------------------------------
\subsection{Three-Way: Workload $\times$ Router $\times$ Pool}
% -----------------------------------------------------------------------------

\begin{opportunity}[Online co-adaptation of compression rate and pool boundary]
\label{opp:workload-shaping}
\Cref{opp:dynamic-boundary} adjusts the pool boundary $\boundary^*$
as the workload CDF shifts, but holds the compression rate $\gamma$
fixed.  FleetOpt~\cite{fleetopt2026} shows that $\gamma$ and
$\boundary^*$ are coupled: the optimal compression rate depends on
the pool boundary, and the optimal pool boundary depends on how
aggressively the router compresses.  Adapting one without the other
leaves cost on the table.

The coupling is workload-dependent.  Archetype~1 (chat-dominated,
concentrated-below) benefits from aggressive compression: most
requests are short, so compressing borderline requests shifts them
into the short pool cheaply.  Archetype~3 (agent-dominated,
concentrated-above) does \emph{not} benefit from aggressive
compression: most requests are long, and compression cannot shrink
them below $\boundary^*$.  As the workload rotates through the
day (Archetype~1 during business hours, Archetype~3 overnight),
the optimal $\gamma$ rotates with it.

The algorithm is a direct extension of \Cref{opp:dynamic-boundary}:
every $T$ minutes, (1)~estimate the current CDF from the
sliding-window sample, (2)~evaluate FleetOpt's closed-form cost
function over a grid of $(\gamma, \boundary^*)$ pairs (the search
space is two-dimensional and small), (3)~apply the minimum-cost
pair by updating the gateway compression threshold and reassigning
pool instances.  Step~(2) reuses FleetOpt's existing solver; the
only new component is the $\gamma$ dimension in the search.

\begin{frisk}
The grid search over $(\gamma, \boundary^*)$ adds negligible
overhead: FleetOpt's cost function evaluates in microseconds, and a
$20 \times 20$ grid covers the practical range.  The pattern of
jointly re-tuning coupled configuration parameters from live
workload telemetry is established: Google
Autopilot~\cite{rzadca2020autopilot} continuously co-adapts CPU
and memory limits for containers, reducing resource slack from
46\% to 23\%, and uses hysteresis to avoid oscillation between
configurations.  The same hysteresis principle applies here: only
apply a new $(\gamma, \boundary^*)$ pair if the cost improvement
exceeds a threshold~$\epsilon$.  The first validation step is to
measure the cost gap between joint $(\gamma, \boundary^*)$
optimization and boundary-only optimization on FleetOpt's three
archetype CDFs.
\end{frisk}
\end{opportunity}

\begin{opportunity}[Mixed-archetype fleet provisioning]
\label{opp:mixed-archetype}
FleetOpt~\cite{fleetopt2026} and FleetSim~\cite{fleetsim2026}
optimize fleet sizing for a single workload CDF archetype.
Production fleets serve mixed workloads simultaneously: a
customer-support chatbot (Archetype~1, concentrated-below), an
internal RAG pipeline (Archetype~2, dispersed), and a coding
agent (Archetype~3, concentrated-above).

The semantic router's domain and workload-type classification
provides the per-request archetype breakdown in real time.  This
breakdown is the missing input to fleet provisioning: a fleet
serving 60\% Archetype~1 and 40\% Archetype~3 has complementary
pool demands (the short pool is busy during business hours when
chat dominates; the long pool is busy overnight when agents run
batch tasks).  Extending FleetOpt to accept a \emph{mixture} of
CDFs---weighted by the router's classification
distribution---would produce fleet sizes that exploit this
complementarity, rather than provisioning for the worst-case
single archetype.

\begin{frisk}
The principle that mixed workloads pack more efficiently than
segregated ones is well established in cluster scheduling.
Borg~\cite{verma2015borg} showed that segregating production
services from batch jobs would require 20--50\% more machines
than running them together, because their resource demands are
complementary (production is memory-heavy and steady; batch is
CPU-heavy and bursty).
Tetris~\cite{grandl2014tetris} formalized this as
multi-dimensional bin packing: by aligning task resource vectors
with available machine capacity across CPU, memory, disk, and
network, it achieved 30--40\% makespan improvements on YARN
clusters with heterogeneous jobs.
DRF~\cite{ghodsi2011drf} provided the fairness foundation,
showing that multi-resource allocation across heterogeneous
demand profiles is both Pareto-efficient and strategy-proof
when each tenant's dominant resource share is equalized.

The LLM inference analogue is direct: archetypes differ in their
dominant resource.  Archetype~1 (short prompts, high request rate)
is prefill-throughput-limited; Archetype~3 (long prompts, low
request rate) is KV-cache-memory-limited.  A fleet provisioned
for the worst-case single archetype over-provisions one resource
and under-provisions the other.  FleetOpt's cost model is
per-archetype; extending it to weighted CDF mixtures is
analytically straightforward (the effective CDF is a mixture
distribution).  The research question is whether the complementarity
is large enough to matter: if the cost function is flat near the
optimum, mixed-archetype provisioning adds complexity without
meaningful savings.  FleetSim~\cite{fleetsim2026} can evaluate this
empirically by simulating mixed workloads across the three archetype
CDFs.  A second risk is classification accuracy: the router's
archetype labels must be reliable, since misclassification shifts
the effective CDF and can worsen fleet sizing.
\end{frisk}
\end{opportunity}

\begin{opportunity}[Energy efficiency as a composable routing objective]
\label{opp:energy-signal}
The 1/W law~\cite{onewlaw2026} provides a closed-form tok/W model
as a function of context length, pool configuration, and GPU
generation.  Two-pool routing recovers most of the
$2$--$4\times$ energy loss caused by the chat-to-agent workload
shift, but only if pool boundaries are energy-aware.  Currently,
the 1/W law is an offline analysis tool; it does not influence
routing decisions at request time.

The opportunity is to operationalize the 1/W model as a composable
signal in the vSR framework, alongside cost, latency, and quality.
For each candidate pool assignment, the router evaluates the
expected tok/W (from the 1/W model given the request's context
length and the pool's GPU type) and incorporates it into the
decision rule.  Operators can then set energy SLOs (``fleet-wide
tok/W $\geq X$'') that the router enforces without sacrificing
cost or latency constraints---or explicitly accept a latency
penalty for energy gains.

\begin{frisk}
The 1/W model is validated and fast to evaluate (closed-form
lookup, sub-millisecond).  Adding it as a signal type in the vSR
framework is engineering work.  Energy-aware and carbon-aware LLM
routing is an emerging research direction:
GreenServ~\cite{greenserv2026} uses multi-armed bandits to route
queries by context-aware features, achieving 31\% energy reduction
with 22\% accuracy improvement over random routing;
GAR~\cite{gar2026} formulates carbon-aware routing as constrained
multi-objective optimization over per-request CO$_2$ estimates
while satisfying accuracy and p95-latency SLOs; and
CarbonFlex~\cite{carbonflex2026} achieves $\sim$57\% carbon
reduction via workload-aware provisioning.  The research question
is multi-objective balancing: how should energy trade off against
cost and latency?  For MoE models, which the 1/W law shows achieve
$5.1\times$ better tok/W~\cite{onewlaw2026}, the energy signal
aligns with cost; for dense models, the objectives may conflict.
Pareto-front analysis across the three CDF archetypes is needed to
characterize when the energy signal changes routing decisions and
when it is dominated by cost.
\end{frisk}
\end{opportunity}

\begin{opportunity}[Governance-as-code for the WRP architecture]
\label{opp:governance}
The \vsr{} DSL~\cite{vllmsr2026} and SIRP
protocol~\cite{sirp2025} provide the building blocks for
policy-as-code.  Extending the DSL to express WRP-level
constraints---``PII queries must stay on-premise,'' ``agent
workloads must route through the high-accuracy pool,'' ``cost
must not exceed \$X/request''---would enable governance-as-code
over the full inference architecture, with compile-time
verification that the constraints are satisfiable.
MPExt~\cite{mpext2025} provides the multi-provider substrate for
cross-fleet enforcement.

\paragraph{RBAC predicates on caller role.}
The bearer token (\Cref{sec:workload:identity}) extends the DSL's
constraint vocabulary from fleet-level properties (pool, cost,
region) to \emph{caller-level} properties (role, scope,
authorization claims).  Concretely:
\texttt{non\_owner.tools.deny~=~[shell, email\_delete,
config\_modify]} expressed in the vSR DSL, validated against the
bearer token's JWT claims at dispatch time.  OPA's Rego language
natively expresses such identity-based policies, and the Agent
Authorization Profile (AAP)~\cite{aap2026} provides the claim
structure---five structured JWT claims covering agent identity,
capabilities with constraints, task binding, delegation tracking, and
audit trails.  Adding RBAC predicates to the governance DSL turns
policy authoring into the foundation for runtime enforcement
(\Cref{opp:rbac-enforcement}).

\begin{frisk}
The DSL and SIRP protocol exist; extending them with pool-level
and cost-level constraints is language design and compiler
engineering.  The policy-as-code paradigm is mature in
infrastructure governance: Open Policy Agent
(OPA)~\cite{openpolicyagent}, a CNCF graduated project, decouples
policy from application logic via a declarative language (Rego) and
serves as Kubernetes admission controller for constraint
enforcement, audit trails, and compile-time satisfiability checks.
The WRP governance layer follows the same separation:
operators author constraints in the vSR DSL, the compiler verifies
satisfiability against the current fleet topology, and the router
enforces them at dispatch time.  Adding RBAC predicates introduces a
new satisfiability dimension: can every authorized role reach at least
one model through at least one pool?  The compiler can verify this
statically against the current fleet topology.  The practical risk is
policy \emph{maintenance}: as the fleet evolves (new GPU types, new
models, new pool topologies), existing constraints may become
unsatisfiable, and operators need clear diagnostics---the same
challenge OPA Gatekeeper addresses with its audit functionality.
\end{frisk}
\end{opportunity}


% -----------------------------------------------------------------------------
\subsection{Capstone: Closed-Loop Self-Adaptation}
% -----------------------------------------------------------------------------

\begin{opportunity}[Closed-loop self-adaptation across WRP knobs]
\label{opp:self-adapt}
The preceding twenty opportunities each adjust one or two WRP
parameters---model selection weights, compression rate~$\gamma$,
pool boundary~$\boundary^*$, retention TTL, cascade threshold,
token budget, joint tool--model selection, safety escalation
threshold, RBAC enforcement rules, caller reputation scores,
behavioral commitment blocks---based on a dedicated feedback signal.  In isolation, each loop is well-defined.  In
combination, the loops interact: adjusting model selection
(\Cref{opp:failure-routing}) changes the session-length
distribution, which shifts the workload CDF, which changes the
optimal pool boundary (\Cref{opp:dynamic-boundary}), which changes
when cascades are worthwhile (\Cref{opp:pool-cascade}).  Tuning
knobs independently ignores these coupling effects---the same
pitfall that database auto-tuning
discovered~\cite{onlinetune2022}: changing one knob invalidates the
optimal setting of another.
\Cref{fig:knob-interactions} maps the dependency chains.

\begin{figure}[t]
\centering
\begin{tikzpicture}[
  knob/.style={rectangle, rounded corners=3pt, draw, minimum height=6mm,
               minimum width=22mm, font=\footnotesize\sffamily,
               inner sep=2pt, align=center},
  wknob/.style={knob, fill=workloadblue!15, draw=workloadblue!60},
  rknob/.style={knob, fill=routergreen!15, draw=routergreen!60},
  pknob/.style={knob, fill=poolorange!15, draw=poolorange!60},
  dep/.style={-{Stealth[length=4pt]}, thick, gray!70},
  bidep/.style={{Stealth[length=4pt]}-{Stealth[length=4pt]}, thick, gray!70},
  every node/.append style={font=\scriptsize},
  node distance=10mm and 14mm,
]
% --- Workload row ---
\node[wknob] (archetype) {Archetype\\composition\\{\tiny Opp\,11,\,13}};
\node[wknob, right=18mm of archetype] (tokenbudget) {Token budget\\{\tiny Opp\,5}};

% --- Router row ---
\node[rknob, below left=12mm and 4mm of archetype] (modelsel)
  {Model selection\\{\tiny Opp\,1,\,2,\,4}};
\node[rknob, right=10mm of modelsel] (gamma)
  {Compression $\gamma$\\{\tiny Opp\,12}};
\node[rknob, right=10mm of gamma] (toolprune)
  {Tool + model\\{\tiny Opp\,17,\,18}};
\node[rknob, right=10mm of toolprune] (cascade)
  {Cascade threshold\\{\tiny Opp\,9}};
\node[rknob, right=10mm of cascade] (retention)
  {Retention TTL\\{\tiny Opp\,10}};

% --- Security row ---
\definecolor{securityred}{RGB}{180,40,60}
\tikzset{sknob/.style={knob, fill=securityred!12, draw=securityred!60}}
\node[sknob, below right=12mm and -2mm of retention] (authz)
  {RBAC + reputation\\{\tiny Opp\,19,\,20,\,21}};

% --- Pool row ---
\node[pknob, below=14mm of gamma] (taustar)
  {Pool boundary $\tau^*$\\{\tiny Opp\,6,\,12}};
\node[pknob, right=14mm of taustar] (poolscale)
  {Per-pool scaling\\{\tiny Opp\,11,\,13}};

% --- Dependency edges ---
\draw[dep] (modelsel) -- node[left, pos=0.4] {\tiny CDF shift} (taustar);
\draw[bidep] (gamma) -- node[right, pos=0.4] {\tiny coupled} (taustar);
\draw[dep] (taustar) -- node[below, pos=0.5] {\tiny load change} (cascade);
\draw[dep] (tokenbudget) -- node[right, pos=0.4]
  {\tiny compress/downgrade} (taustar);
\draw[dep] (archetype) -- node[left, pos=0.4]
  {\tiny demand forecast} (poolscale);
\draw[dep] (poolscale) -- node[below, pos=0.5]
  {\tiny pool pressure} (cascade);
\draw[dep] (retention) -- node[right, pos=0.4]
  {\tiny cache hit rate} (cascade);
\draw[dep] (modelsel.south east) to[bend right=15]
  node[below, pos=0.5] {\tiny failure rates} (modelsel.south);
\draw[dep, routergreen!60] (modelsel) to[bend left=20]
  node[above, pos=0.5] {\tiny context pollution} (tokenbudget);
\draw[dep] (toolprune) -- node[above, pos=0.5]
  {\tiny context savings} (tokenbudget);
\draw[dep, securityred!60] (authz) to[bend left=25]
  node[below, pos=0.5] {\tiny tool restrict} (toolprune);
\draw[dep, securityred!60] (authz) to[bend right=20]
  node[right, pos=0.4] {\tiny budget cap} (tokenbudget);

% --- Dimension labels ---
\node[above=2mm of archetype, font=\footnotesize\sffamily\bfseries,
      color=workloadblue] {Workload};
\node[left=2mm of modelsel, font=\footnotesize\sffamily\bfseries,
      color=routergreen] {Router};
\node[below left=2mm and -4mm of taustar,
      font=\footnotesize\sffamily\bfseries,
      color=poolorange] {Pool};
\node[right=2mm of authz, font=\footnotesize\sffamily\bfseries,
      color=securityred] {Security};
\end{tikzpicture}
\caption{Knob-interaction dependencies across WRP opportunities.
Adjusting one knob (e.g., model selection) cascades through the
workload CDF into pool-layer parameters ($\tau^*$, cascade threshold),
requiring coordinated adaptation.  The green arrow from model
selection to token budget captures the memory-management feedback
loop: poor model choices cause failures that pollute session context,
driving token-budget explosion.  Bidirectional arrow indicates
the $\gamma$--$\tau^*$ coupling from \Cref{opp:workload-shaping}.}
\label{fig:knob-interactions}
\end{figure}

The opportunity is a \emph{meta-controller} that orchestrates the
individual adaptation loops as a unified system.  The architecture
follows the Monitor--Analyze--Plan--Execute (MAPE-K) reference
model from autonomic computing~\cite{kephart2003autonomic}:
\begin{enumerate}[nosep]
  \item \textbf{Monitor.}  Collect a per-session outcome record
        that unifies all signals: model sequence, tool outcomes,
        hallucination verdicts (HaluGate), safety scores
        (SafetyL1/L2), quality signal (Feedback Detector), turn
        count, token consumption, TTFT, context growth rate,
        failure-driven token waste (tokens consumed by turns that
        produced errors subsequently corrected or discarded), and
        task-completion indicator.  The cross-session failure
        memory---aggregated per (model, domain, failure
        class)---provides the fleet-wide baseline against which
        individual session outcomes are compared.
  \item \textbf{Analyze.}  Map each poor outcome to a failure class
        \emph{and} the knob that could have prevented it:
        ``12-turn session $\to$ wrong model at turn~2
        (\Cref{opp:session-cascade}),'' ``SLO miss $\to$ strong
        pool overloaded at cascade (\Cref{opp:pool-cascade}),''
        ``token explosion $\to$ no budget enforcement
        (\Cref{opp:token-budget}),'' ``40\% token waste $\to$
        hallucination at turn~3 polluted context for turns 4--8,
        domain-specific compression would have pruned the error
        trace (\Cref{opp:token-budget}).''  Cross-session
        failure schemata~\cite{correct2025} accelerate this
        analysis: structurally similar failures recur across
        sessions, so a new session's failure can be matched to a
        known schema rather than diagnosed from scratch.
  \item \textbf{Plan.}  Select knob adjustments that address the
        dominant failure classes while respecting interaction
        constraints: e.g., if model-selection changes shift the CDF
        enough to invalidate the current $\boundary^*$, schedule a
        joint $(\gamma, \boundary^*)$ re-optimization
        (\Cref{opp:workload-shaping}) in the same cycle.
  \item \textbf{Execute.}  Apply the new configuration via staged
        rollout: evaluate the candidate policy offline via
        importance-weighted off-policy
        estimation~\cite{deltope2024}, then route 5\% canary
        traffic under the new policy, measure the composite
        outcome, and promote to full traffic only if the outcome
        improves beyond a hysteresis threshold.
\end{enumerate}

\Cref{tab:mape} maps each opportunity to its MAPE-K role: the
observable it monitors, the knob it tunes, and the metric it
evaluates improvement against.

\begin{table}[t]
\centering
\caption{MAPE-K role of each opportunity: what the adaptation loop
monitors, which knob it adjusts, and how it evaluates improvement.}
\label{tab:mape}
\scriptsize
\renewcommand{\arraystretch}{1.10}
\begin{tabularx}{\textwidth}{c l X X X}
\toprule
\textbf{\#} & \textbf{Opportunity} & \textbf{Monitor (observable)}
  & \textbf{Knob (control)} & \textbf{Evaluate (metric)} \\
\midrule
\rowcolor{routergreen!6}
1  & Session-length RL
   & Session trajectories (model, tool, outcome per turn)
   & Model + tool selection policy
   & Median session length (turns, tokens) \\
2  & HaluGate reputation
   & Per-response hallucination verdicts
   & Per-(model, domain) routing weights
   & Hallucination rate per domain \\
\rowcolor{routergreen!6}
3  & Cumulative risk
   & Per-turn SafetyL1 confidence scores
   & Safety escalation threshold
   & FPR vs.\ attack catch rate \\
4  & Failure-class routing
   & Multi-signal outcome per (model, domain, failure class)
   & Failure-class routing weights
   & Oracle-utility capture rate \\
\rowcolor{routergreen!6}
5  & Token budget
   & Per-session token counter + failure-driven context bloat
   & Budget threshold + domain-specific compression policy
   & Token waste from failed sessions \\
\rowcolor{routergreen!6}
17 & Joint tool--model optimization
   & Per-(domain, turn, tool, model) outcomes + transition graph
   & Tool subset + resolving model + transition priors per dispatch
   & Aggregate session cost $\times$ task-completion rate \\
\rowcolor{routergreen!6}
18 & Gateway agent loops
   & Session depth, tool/cache/KV telemetry, cross-tenant outcomes
   & Instruction/tool surface + cache-block order + TTL/admission hints
   & A/B: tokens/rounds/TTFT vs.\ ITR-only or scheduler-only baselines \\
19 & RBAC enforcement
   & Bearer token JWT claims + model tool\_call outputs
   & Per-(role, tool, operation) permission matrix
   & Unauthorized tool-call block rate + false-denial rate \\
\rowcolor{routergreen!6}
20 & Caller reputation
   & Per-bearer-token security events across fleet
   & Reputation score $\to$ tool/model/budget restrictions
   & Cross-agent attack propagation prevention rate \\
21 & Behavioral commitments
   & Model boundary declarations in output text
   & Per-(token, agent) persistent block rules
   & Re-engagement-under-pressure prevention rate \\
\midrule
\rowcolor{poolorange!6}
6  & Follow-the-sun
   & Sliding-window prompt-length CDF
   & Pool boundary $\tau^*$
   & Fleet cost at SLO target \\
7  & Pool-state inference
   & Dispatch/completion counts, TTFT
   & Per-instance routing weights
   & Per-pool SLO attainment \\
\rowcolor{poolorange!6}
8  & Output-length routing
   & Actual completion lengths
   & Pool routing (prompt $+$ predicted output)
   & Short-pool overflow rate \\
9  & Pool-aware cascading
   & Estimated queue depth per pool
   & Cascade-abort threshold
   & Quality vs.\ TTFT SLO \\
\rowcolor{poolorange!6}
10 & KV-cache retention
   & Domain, turn number, OATS tool latency
   & Retention TTL + eviction priority
   & Cache hit rate on agent resume \\
\midrule
11 & Proactive scaling
   & Per-archetype traffic time series
   & Per-pool instance count
   & Forecast accuracy vs.\ aggregate \\
\rowcolor{gray!6}
12 & $(\gamma, \tau^*)$ co-adapt
   & Sliding-window CDF + current $(\gamma, \tau^*)$
   & Joint $(\gamma, \tau^*)$
   & Cost gap vs.\ boundary-only \\
13 & Mixed-archetype fleet
   & Router archetype breakdown
   & Fleet sizing weights
   & Complementarity savings \\
\rowcolor{gray!6}
14 & Energy routing
   & tok/W per pool configuration
   & Energy weight in decision rule
   & Fleet-wide tok/W at SLO \\
15 & Governance-as-code
   & Policy constraints + fleet topology
   & DSL rule set
   & Constraint satisfiability \\
\rowcolor{gray!6}
16 & Self-adaptation
   & All of the above (composite session + security record)
   & All knobs jointly (incl.\ RBAC, reputation, commitments)
   & Composite session outcome + security posture \\
\bottomrule
\end{tabularx}
\end{table}

The composite outcome function is the reward signal that ties the
system together.  It must weight multiple objectives---session
length, token cost, SLO attainment, failure rate, safety---into a
single scalar.  Recent work shows this is tractable:
PROTEUS~\cite{proteus2026} uses Lagrangian RL to accept accuracy
targets as runtime input without retraining;
BaRP~\cite{barp2025} trains under bandit feedback with
preference-tunable cost-quality tradeoffs, outperforming offline
routers by 12.5\%.  The meta-controller extends these ideas from
model selection to the full WRP control surface.

The router's structural advantage is again fleet-wide visibility:
it observes all sessions, all outcomes, and all knob settings
simultaneously.  No individual agent, model provider, or pool
instance has this view.  The meta-controller is the natural
convergence point where the individual adaptation loops become a
self-improving system.

\begin{frisk}
The MAPE-K pattern is mature in autonomic
computing~\cite{kephart2003autonomic} and has been applied to
database configuration tuning, where OnlineTune~\cite{onlinetune2022}
demonstrated safe online multi-knob adaptation in production cloud
databases (14--165\% improvement, 91--99.5\% fewer unsafe
recommendations) by embedding workload context and using contextual
Bayesian optimization with safe exploration.  The LLM inference
analogue replaces database knobs with WRP parameters and replaces
query throughput with the composite session-outcome metric.  The
research questions are: (a)~dimensionality---how many WRP knobs
can be jointly tuned before the search space becomes intractable
(OnlineTune handles $\sim$50 database knobs; WRP has $\sim$10,
well within range); (b)~evaluation cost---off-policy evaluation
from logged sessions avoids the need for live A/B tests on every
candidate configuration; (c)~interaction
strength---if the coupling between knobs is weak (e.g., changing
model selection does not materially shift the CDF), independent
loops suffice and the meta-controller adds unnecessary complexity.
FleetSim~\cite{fleetsim2026} can quantify the interaction strength
empirically.  The baseline is independent per-knob adaptation; the
metric is composite session outcome (quality-weighted session cost)
on a multi-week FleetSim trace with diurnal workload rotation.
\end{frisk}
\end{opportunity}


